% =============================================================================
%  acl/main.tex -- MINIMAL ACL / NeurIPS-workshop style stub
%
%  This is a stub, not the official template. It uses the plain `article` class
%  plus natbib -- which is exactly what both official style packages assume --
%  so it compiles today, on a stock TeX Live, with no download. Turning it into
%  a real submission is a two-line change (see the switch block below).
%
%  BEFORE SUBMISSION:
%    1. download the official style files yourself:
%         ACL     https://github.com/acl-org/acl-style-files            (verify)
%         NeurIPS https://neurips.cc  ->  Paper Information -> Style Files (verify)
%    2. put acl.sty + acl_natbib.bst (or neurips_20XX.sty) in THIS directory,
%       read the LICENCE file that ships with them, and keep that licence file
%       here. Do not redistribute the official files from this repo.
%    3. read the current page limit and the mandatory checklist (ACL:
%       Responsible NLP checklist, question E1 covers AI use; NeurIPS: paper
%       checklist, question on LLM use).
%
%  Compile from THIS directory (relative \input paths depend on it):
%      cd 05_paper/template/acl && latexmk -pdf main.tex
% =============================================================================

\documentclass[11pt]{article}

% --- SWITCH TO THE OFFICIAL STYLE --------------------------------------------
% ACL (review mode; use \usepackage{acl} for camera-ready):
%     \usepackage[review]{acl}
%     \bibliographystyle{acl_natbib}
% NeurIPS workshop / main track (pick the file name that matches the year):
%     \usepackage[preprint]{neurips_2026}     % or [final]
%     \bibliographystyle{plainnat}            % neurips ships no .bst of its own
% Keep the natbib load below until you switch; acl.sty loads natbib itself and
% will error on a duplicate load, so delete it when you switch.
% -----------------------------------------------------------------------------

\usepackage[utf8]{inputenc}
\usepackage[T1]{fontenc}
\usepackage[margin=1in]{geometry}
\usepackage{amsmath,amssymb}
\usepackage{graphicx}
\usepackage{booktabs}
\usepackage[round,authoryear]{natbib}
\usepackage{xcolor}
\usepackage{url}
\usepackage[hidelinks]{hyperref}

\newcommand{\method}{TODO-Method-Name}
\newcommand{\fillme}[1]{%
  \par\noindent\fbox{\parbox{0.96\linewidth}{\small\textbf{TODO} --- #1}}\par}

\title{TODO(title): specific and falsifiable\\
  \large TODO(short title for the running head, $\leq$ 8 words)}

\author{%
  TODO(Author One)\\
  TODO(Affiliation)\\
  \texttt{TODO(email)}
  \and
  TODO(Author Two)\\
  TODO(Affiliation)\\
  \texttt{TODO(email)}
  \thanks{%
    \textbf{Use of AI tools.} TODO(disclosure): name each tool and version, the
    sections it touched, and what it did (e.g.\ ``Section~3 was written with
    inputs from <tool>, then rewritten by the author''). ACL requires AI use to
    be fully disclosed; TODO(venue): quote the policy URL and the retrieval
    date, and check whether the venue wants this in the Acknowledgements
    instead. No AI tool produced data, results, figures or references; every
    citation is recorded in \texttt{02\_literature/verified\_citations.csv}.
    The tools are not authors. See Section~\ref{sec:ethics}.}%
}

\date{}

\begin{document}
\maketitle

\begin{abstract}
% TODO(abstract): 150-200 words for ACL, <= 250 for NeurIPS. Problem, method,
% headline number (% num:<id> marker), limitation. No citations.
\fillme{Abstract: problem / method / headline result with its index id /
limitation.}
\end{abstract}

% -----------------------------------------------------------------------------
\section{Introduction}
\label{sec:intro}
% TODO(intro): problem / gap / contributions / headline result, and the AI
% disclosure obligation: ACL requires AI use to be disclosed in the
% Acknowledgements (see Section~\ref{sec:ethics} and the title footnote).
\fillme{Introduction: problem, gap, contributions, headline result.}

\section{Related Work}
\label{sec:related}
% TODO(related): clusters of prior work, then the delta. Verified keys only.
\fillme{Related Work: clusters + delta.}

\section{Method}
\label{sec:method}
% TODO(method): notation, method, cost.
\fillme{Method: notation, method, cost.}

\section{Experiments}
\label{sec:exp}
% TODO(exp): data, baselines, protocol (seeds, runs, command), metrics.
\fillme{Experiments: data, baselines, protocol, metrics.}

\section{Results}
\label{sec:results}
% TODO(results): the answer first, then Table~\ref{tab:main-results} and
% Figure~\ref{fig:main-accuracy}. Numbers are generated, never typed.
\fillme{Results: the answer, then the evidence, then what did not work.}

\begin{table}[t]
  \centering
  \caption{TODO(caption): what is compared, on what data, over how many seeds.}
  \label{tab:main-results}
  \IfFileExists{../tables/generated/main-results.tex}{%
    \input{../tables/generated/main-results.tex}%
  }{%
    \fbox{\parbox{0.9\linewidth}{\small No generated table yet: run
      \texttt{tables/make\_tables.py} (see BUILD.md).}}%
  }
\end{table}

\begin{figure}[t]
  \centering
  \IfFileExists{../figures/generated/main-accuracy.pdf}{%
    \includegraphics[width=0.85\linewidth]{../figures/generated/main-accuracy.pdf}%
  }{%
    \fbox{\parbox[c][3.4cm][c]{0.8\linewidth}{\centering\small Figure PDF not
      built (needs matplotlib): run \texttt{figures/make\_figures.py}.}}%
  }
  \caption{TODO(caption): axes, unit, uncertainty, and the one thing to notice.}
  \label{fig:main-accuracy}
\end{figure}

\section{Limitations}
\label{sec:limitations}
% TODO(limitations): ACL has a dedicated Limitations section and reviewers
% score it. Data, compute, statistics, scope -- and what would fix each.
\fillme{Limitations: data, compute, statistics, scope.}

\section{Ethics Statement}
\label{sec:ethics}
% TODO(ethics): data licences, annotator compensation if any, misuse, and the
% AI-assistance statement. ACL wants it "fully disclosed", e.g. "Section 3 was
% written with inputs from <tool>". Same wording as the footnote.
\fillme{Ethics: licences, misuse, AI-assistance statement.}

\section{Reproducibility Statement}
\label{sec:repro}
% TODO(repro): repo, command, seeds, released artefacts, rerun tolerance.
\fillme{Reproducibility: repo, command, seeds, tolerance.}

\section{Conclusion}
\label{sec:conclusion}
% TODO(conclusion): what is now known, calibrated to the evidence.
\fillme{Conclusion: what is now known, and the next experiment.}

% -----------------------------------------------------------------------------
\section*{Acknowledgements}
% TODO(ack): non-author help, funding, and -- for ACL -- the AI-use statement
% if the venue asks for it here rather than in a footnote. Keep one wording.
\fillme{Acknowledgements: help, funding, AI-use statement.}

\bibliographystyle{plainnat}
\bibliography{../refs}

\end{document}
